% =============================================================
% STANDALONE COMPILABLE LaTeX DOCUMENT
% NDVI Novel Methodology — Forest Fire Susceptibility Mapping
% IIT Roorkee | Dept. Applied Mathematics & Scientific Computing
% Supervisor: Prof. Millie Pant
%
% HOW TO COMPILE IN TeXstudio:
%   1. Set Bibliography: Options → Configure TeXstudio →
%      Build → Default Bibliography Tool → BibTeX
%   2. Press F5 (Build & View) OR run in order:
%      pdflatex → bibtex → pdflatex → pdflatex
%
% COMPILER: pdflatex (NOT XeLaTeX or LuaLaTeX)
% =============================================================

\documentclass[12pt, a4paper]{article}

% ── Essential Packages ──────────────────────────────────────
\usepackage[utf8]{inputenc}
\usepackage[T1]{fontenc}
\usepackage[english]{babel}
\usepackage{geometry}
\geometry{
    a4paper,
    left   = 2.5cm,
    right  = 2.5cm,
    top    = 2.5cm,
    bottom = 2.5cm
}

% ── Mathematics ─────────────────────────────────────────────
\usepackage{amsmath}
\usepackage{amssymb}
\usepackage{bm}

% ── Graphics & Tables ───────────────────────────────────────
\usepackage{graphicx}
\usepackage{subcaption}
\usepackage{booktabs}
\usepackage{array}
\usepackage{multirow}
\usepackage{longtable}

% ── Algorithms ──────────────────────────────────────────────
\usepackage{algorithm}
\usepackage{algpseudocode}

% ── References (natbib — required for \citet, \citep) ───────
\usepackage[round, authoryear, sort]{natbib}

% ── Colours & Hyperlinks ────────────────────────────────────
\usepackage{xcolor}
\usepackage[
    colorlinks = true,
    linkcolor  = blue,
    citecolor  = teal,
    urlcolor   = blue
]{hyperref}

% ── Miscellaneous ───────────────────────────────────────────
\usepackage{setspace}
\onehalfspacing
\usepackage{parskip}
\setlength{\parindent}{0pt}
\usepackage{enumitem}
\usepackage{caption}
\captionsetup{font=small, labelfont=bf}
\usepackage{float}

% ── Custom Colours ──────────────────────────────────────────
\definecolor{novelblue}{RGB}{0, 82, 147}
\definecolor{novelbg}{RGB}{240, 248, 255}

% ── Custom Environments ─────────────────────────────────────
\usepackage{tcolorbox}
\tcbuselibrary{skins, breakable}
\newtcolorbox{novelbox}[1]{
    colback   = novelbg,
    colframe  = novelblue,
    fonttitle = \bfseries,
    title     = {#1},
    breakable = true
}

% ── Header & Footer ─────────────────────────────────────────
\usepackage{fancyhdr}
\pagestyle{fancy}
\fancyhf{}
\rhead{\small NDVI Methodology | IIT Roorkee}
\lhead{\small Forest Fire Susceptibility Mapping — India}
\cfoot{\thepage}

% =============================================================
\begin{document}
% =============================================================

% ── Title Page ──────────────────────────────────────────────
\begin{titlepage}
\centering
\vspace*{2cm}

{\LARGE \bfseries Novel NDVI-Based Vegetation Stress Analysis\\[0.5em]
for Forest Fire Susceptibility Mapping in India}

\vspace{1.5cm}

{\large \textit{Methodology Section — Research Working Document}}

\vspace{1cm}

\rule{0.8\linewidth}{0.4pt}

\vspace{1cm}

{\large \textbf{[Your Name]}}\\[0.3em]
{\normalsize Senior Research Fellow (SRF)}\\[0.2em]
{\normalsize Department of Applied Mathematics \& Scientific Computing}\\[0.2em]
{\normalsize Indian Institute of Technology Roorkee}

\vspace{0.5cm}

{\normalsize Under the supervision of \textbf{Prof. Millie Pant}}

\vspace{1cm}

\rule{0.8\linewidth}{0.4pt}

\vspace{0.8cm}

{\normalsize \textbf{Data:} MOD13C2 / MOD13A3 v061, NASA LP DAAC, 2000--2022}\\[0.3em]
{\normalsize \textbf{Extends:} \citealt{biswas2025} —
             \textit{Environ.\ Sci.\ Pollut.\ Res.}, 32:4856--4878}\\[0.3em]
{\normalsize \textbf{Funding:} Ministry of Education, Government of India}

\vspace{1.5cm}

{\normalsize \today}

\end{titlepage}

% ── Table of Contents ───────────────────────────────────────
\tableofcontents
\newpage

% =============================================================
\section{NDVI-Based Vegetation Stress Analysis for Forest Fire
         Susceptibility}
\label{sec:ndvi_methodology}
% =============================================================

\subsection{Overview and Motivation}
\label{sec:overview}

The Normalized Difference Vegetation Index (NDVI) is defined as:

\begin{equation}
    \text{NDVI} = \frac{\rho_{\text{NIR}} - \rho_{\text{Red}}}
                       {\rho_{\text{NIR}} + \rho_{\text{Red}}}
    \label{eq:ndvi_definition}
\end{equation}

where $\rho_{\text{NIR}}$ and $\rho_{\text{Red}}$ denote the surface
reflectance in the Near-Infrared and Red spectral bands respectively.
The index ranges from $-0.2$ to $1.0$, with higher values indicating
denser and healthier vegetation.

NDVI emerged as the \textbf{single most important predictor} of forest
fire occurrence in India among 15 environmental variables, contributing
$22.3\%$ permutation importance and $28.4\%$ percentage contribution
in the MaxEnt framework of \citet{biswas2025}. Despite this primacy,
prior studies — including \citet{biswas2025} — employed NDVI solely as a
static or temporally-averaged predictor, foregoing the rich temporal
information encoded in the 22-year MODIS satellite record (2000--2022).

\begin{novelbox}{Methodological Gap Addressed in This Study}
\citet{biswas2025} applied NDVI as:
\begin{itemize}[leftmargin=1.5em]
    \item A single time-averaged raster layer
    \item Without pixel-level quality filtering (cloud/snow contamination retained)
    \item Without temporal decomposition (seasonal cycle confounds prediction)
    \item Without pre-fire stress signal extraction
    \item Assuming a linear NDVI--fire relationship
\end{itemize}
This study introduces five mathematically grounded extensions that
transform raw NDVI into physically interpretable, temporally dynamic,
and spatially aware predictors of fire susceptibility.
\end{novelbox}

The five novel components introduced in this study are:

\begin{enumerate}
    \item \textbf{QA-filtered scale-corrected NDVI} (Section~\ref{sec:ndvi_data})
    \item \textbf{Seasonal decomposition and anomaly extraction}
          (Section~\ref{sec:ndvi_decomp})
    \item \textbf{Pre-fire Cumulative Vegetation Stress Index (CVSI)}
          (Section~\ref{sec:cvsi})
    \item \textbf{Spatial autocorrelation and fire-cluster alignment}
          via Moran's~$I$ (Section~\ref{sec:morans_i})
    \item \textbf{NDVI--fire threshold detection via breakpoint analysis}
          (Section~\ref{sec:threshold})
\end{enumerate}

% =============================================================
\subsection{Data Source and Preprocessing}
\label{sec:ndvi_data}
% =============================================================

NDVI data were obtained from the MODIS/Terra Vegetation Indices Monthly
product \textbf{MOD13C2 Version~6.1} \citep{didan2021}, hosted at the
NASA Land Processes Distributed Active Archive Center (LP~DAAC) and
accessed via the AppEEARS extraction service \citep{appeears2023}.
The product provides monthly composite NDVI at $0.05^{\circ}$
(${\approx}5.6$~km) spatial resolution in the Climate Modeling Grid
(CMG) projection. A higher-resolution variant (MOD13A3 v061, 1~km
monthly) was also evaluated and yields improved spatial detail over the
$0.05^{\circ}$ product used in \citet{biswas2025}. The study period
spans February 2000 through August 2022 ($T = 272$ monthly composite
images).

% ── 2.1 ─────────────────────────────────────────────────────
\subsubsection{Scale Correction and Fill-Value Masking}
\label{sec:scale}

Raw MOD13C2 stores NDVI as signed 16-bit integers. The corrected
floating-point NDVI for pixel $(i,j)$ at time step $t$ is:

\begin{equation}
    \text{NDVI}_{ij}^{(t)} =
    \begin{cases}
        \text{raw}_{ij}^{(t)} \times 10^{-4} &
            \text{if } -2000 \leq \text{raw}_{ij}^{(t)} \leq 10000 \\[4pt]
        \text{NaN} & \text{otherwise}
    \end{cases}
    \label{eq:ndvi_scale}
\end{equation}

yielding true NDVI $\in [-0.2,\, 1.0]$. The designated fill value of
$-3000$ and all values outside the valid range are replaced with
\texttt{NaN} prior to any analysis.

% ── 2.2 ─────────────────────────────────────────────────────
\subsubsection{Quality Assurance Filtering}
\label{sec:qa_filter}

\begin{novelbox}{Novel Contribution~1: Pixel-Level QA Filtering}
\citet{biswas2025} applied no pixel-level quality screening to NDVI,
implicitly including cloud-contaminated and snow/ice-affected retrievals.
This work applies the \texttt{pixel\_reliability} Science Data Set (SDS)
from MOD13C2 to filter each pixel before analysis.
\end{novelbox}

The QA-filtered NDVI is defined as:

\begin{equation}
    \text{NDVI}_{ij}^{(t),\,\text{QA}} =
    \begin{cases}
        \text{NDVI}_{ij}^{(t)} & \text{if } q_{ij}^{(t)} \in \{0,\, 1\} \\[4pt]
        \text{NaN}             & \text{otherwise}
    \end{cases}
    \label{eq:qa_filter}
\end{equation}

where pixel reliability values are defined as:
$q = 0$ (Good data, use without further screening),
$q = 1$ (Marginal data, useful but inspect detailed QA),
$q = 2$ (Snow/Ice — exclude),
$q = 3$ (Cloudy — exclude), and
$q = -1$ (Fill/No data) \citep{didan2021_ug}.
QA filtering is especially critical for monsoon months (June--September),
when cloud contamination over India is highest and unfiltered NDVI values
are systematically negatively biased.

% =============================================================
\subsection{Seasonal Decomposition and Anomaly Extraction}
\label{sec:ndvi_decomp}
% =============================================================

\subsubsection{Climatological Mean and Monthly Anomaly}
\label{sec:anomaly}

India's vegetation phenology is dominated by a strong annual monsoon
cycle. To isolate the inter-annual fire-relevant signal from the
seasonal background, the long-term climatological mean for calendar
month $m$ at pixel $(i,j)$ is computed over the 2001--2020 baseline
(consistent with \citealt{biswas2025}):

\begin{equation}
    \bar{\mu}_{ij}^{(m)} =
    \frac{1}{\lvert\mathcal{Y}_m\rvert}
    \sum_{y \in \mathcal{Y}_m}
    \text{NDVI}_{ij}^{(y,m),\,\text{QA}}
    \label{eq:clim_mean}
\end{equation}

where $\mathcal{Y}_m = \{2001, 2002, \ldots, 2020\}$ and only
non-\texttt{NaN} values contribute. The \textbf{monthly NDVI anomaly}
at year $y$, month $m$, pixel $(i,j)$ is:

\begin{equation}
    \boxed{
    \delta_{ij}^{(y,m)} =
    \text{NDVI}_{ij}^{(y,m),\,\text{QA}} - \bar{\mu}_{ij}^{(m)}
    }
    \label{eq:ndvi_anomaly}
\end{equation}

Negative anomaly ($\delta < 0$) indicates below-average vegetation
density (stress); positive ($\delta > 0$) indicates above-average
greenness. This operation removes the confounding seasonal cycle and
exposes the inter-annual vegetation stress signal directly relevant
to fire ignition.

\subsubsection{STL Decomposition for Trend Isolation}
\label{sec:stl}

\begin{novelbox}{Novel Contribution~2: STL Temporal Decomposition}
Biswas et al.\ \citeyearpar{biswas2025} did not perform any temporal
decomposition of NDVI. This study decomposes the 22-year pixel-wise
time series using STL (Seasonal and Trend decomposition using LOESS),
separating long-term vegetation trend from the seasonal and residual
components.
\end{novelbox}

The STL decomposition \citep{cleveland1990stl} is applied to each
spatial pixel's complete time series:

\begin{equation}
    \text{NDVI}_{ij}^{(t)} = T_{ij}^{(t)} + S_{ij}^{(t)} + R_{ij}^{(t)}
    \label{eq:stl}
\end{equation}

where:
\begin{itemize}
    \item $T_{ij}^{(t)}$ — \textbf{Trend component}: captures inter-annual
          vegetation change driven by deforestation, afforestation, and
          land-use change. A declining trend ($\partial T / \partial t < 0$)
          over a forest pixel indicates vegetation loss and elevated fire risk.
    \item $S_{ij}^{(t)}$ — \textbf{Seasonal component}: the repeating
          annual monsoon-driven phenology cycle (period $= 12$ months).
    \item $R_{ij}^{(t)}$ — \textbf{Residual component}: captures anomalous
          events including drought episodes, pest outbreaks, pre-fire fuel
          drying, and post-fire regeneration.
\end{itemize}

The \textbf{Mann--Kendall trend test} \citep{mann1945, kendall1975} is
applied pixel-wise to $T_{ij}^{(t)}$ to identify statistically significant
greening or browning trends at the $\alpha = 0.05$ significance level.
The Mann--Kendall statistic $\tau$ is:

\begin{equation}
    \tau = \frac{P - Q}{\frac{1}{2}n(n-1)}
    \label{eq:mk_tau}
\end{equation}

where $P$ is the number of concordant pairs and $Q$ the number of
discordant pairs in the time series. Pixels with $\tau < 0$ and
$p < 0.05$ are classified as \textbf{significant browning} zones —
regions of long-term vegetation decline contributing to elevated fire
susceptibility.

% =============================================================
\subsection{Cumulative Vegetation Stress Index (CVSI)}
\label{sec:cvsi}
% =============================================================

\begin{novelbox}{Novel Contribution~3: Pre-Fire Vegetation Stress Quantification}
No prior India-scale forest fire susceptibility study has introduced a
temporally-accumulating pre-fire NDVI stress metric. The CVSI defined
below provides a direct proxy for fuel moisture deficit in the weeks
to months preceding ignition, grounded in combustion physics.
\end{novelbox}

\subsubsection{Mathematical Formulation}

Let $\tau_f$ denote the month of fire occurrence at pixel $(i,j)$.
The $k$-month \textbf{Cumulative Vegetation Stress Index (CVSI)}
leading up to $\tau_f$ is:

\begin{equation}
    \boxed{
    \text{CVSI}_{ij}(\tau_f,\, k) =
    \sum_{t = \tau_f - k}^{\tau_f - 1}
    \max\!\left(-\delta_{ij}^{(t)},\; 0\right)
    }
    \label{eq:cvsi}
\end{equation}

The operator $\max(-\delta, 0)$ retains only negative anomaly episodes
(stress) and discards above-average greenness months. For $k = 3$:

\begin{align}
    \text{CVSI}_{ij}(\tau_f,\, 3) \;=\;&
    \max\!\left(-\delta_{ij}^{(\tau_f - 3)},\, 0\right) \nonumber \\
    +\;& \max\!\left(-\delta_{ij}^{(\tau_f - 2)},\, 0\right) \nonumber \\
    +\;& \max\!\left(-\delta_{ij}^{(\tau_f - 1)},\, 0\right)
    \label{eq:cvsi3}
\end{align}

\textbf{Physical justification:} Prolonged below-average NDVI signals
reduced canopy moisture content, accelerated fine fuel drying, and
increased bulk density of dead surface fuels susceptible to ignition
\citep{chuvieco2004}. The CVSI thus quantifies the \emph{pre-fire
fuel moisture deficit}, a well-established physical precursor to
wildfire ignition \citep{flannigan2000}.

\subsubsection{Optimal Lag Selection via Mutual Information}
\label{sec:mi_lag}

The optimal accumulation window $k^{*}$ is selected by maximising the
\textbf{mutual information} (MI) between $\text{CVSI}(\tau_f, k)$
and binary fire occurrence labels $F_{ij} \in \{0, 1\}$:

\begin{equation}
    k^{*} = \underset{k \;\in\; \{1,2,3,4,5,6\}}{\arg\max}\;
    I\!\left(\text{CVSI}(\tau_f, k)\;;\; F_{ij}\right)
    \label{eq:mi_lag}
\end{equation}

where the Shannon mutual information is:

\begin{equation}
    I(X;\,Y) = \sum_{x,\,y} p(x,y)\,\log_2
    \frac{p(x,y)}{p(x)\,p(y)}
    \label{eq:mutual_info}
\end{equation}

\subsubsection{Comparison with Existing Fire Danger Indices}

Table~\ref{tab:cvsi_comparison} compares the CVSI with commonly used
fire danger indices.

\begin{table}[H]
\centering
\caption{Comparison of CVSI with existing vegetation-based fire indices}
\label{tab:cvsi_comparison}
\begin{tabular}{@{}llll@{}}
\toprule
\textbf{Index} & \textbf{Input} & \textbf{Temporal depth} &
\textbf{Available for India?} \\
\midrule
Raw NDVI \citep{biswas2025}   & NDVI              & Single month      & Yes \\
NDWI                           & NIR, SWIR         & Single month      & Yes \\
FWI (Canadian)                 & Climate variables & Daily rolling     & Limited \\
KBDI                           & Precipitation, T  & Monthly           & Limited \\
\textbf{CVSI (this work)}      & \textbf{NDVI anomaly} &
    \textbf{$k^*$ months pre-fire} & \textbf{Yes (national)} \\
\bottomrule
\end{tabular}
\end{table}

% =============================================================
\subsection{Spatial Autocorrelation: Global and Local Moran's $I$}
\label{sec:morans_i}
% =============================================================

\begin{novelbox}{Novel Contribution~4: Spatial Neighbourhood of Vegetation}
\citet{biswas2025} treated all spatial pixels as independent
observations. This work explicitly models the spatial dependency
structure of NDVI using Moran's~$I$, capturing the fire-relevant
information embedded in a pixel's vegetation neighbourhood.
\end{novelbox}

\subsubsection{Global Moran's $I$}
\label{sec:global_moran}

For a spatial field $\{x_{ij}\}$ with $N$ pixels and row-standardised
spatial weight matrix $\mathbf{W}$ (Queen contiguity):

\begin{equation}
    I = \frac{N}{\displaystyle\sum_{i,j}\sum_{k,l} w_{ij,kl}}
    \cdot
    \frac{\displaystyle\sum_{i,j}\sum_{k,l}
          w_{ij,kl}
          \left(x_{ij} - \bar{x}\right)
          \left(x_{kl} - \bar{x}\right)}
         {\displaystyle\sum_{i,j}\left(x_{ij} - \bar{x}\right)^{2}}
    \label{eq:global_morans_i}
\end{equation}

Under spatial randomness, $\mathbb{E}[I] = -1/(N-1)$. Statistical
significance is assessed via permutation testing ($n_{\text{perm}} = 999$)
at $\alpha = 0.05$. $I > 0$ indicates spatial clustering; $I < 0$
indicates dispersion.

\subsubsection{Local Indicators of Spatial Association (LISA)}
\label{sec:lisa}

The Local Moran's $I$ statistic at pixel $(i,j)$ is:

\begin{equation}
    I_{ij} =
    \frac{x_{ij} - \bar{x}}{S^{2}}
    \sum_{k,l} w_{ij,kl}\left(x_{kl} - \bar{x}\right)
    \label{eq:local_morans_i}
\end{equation}

where $S^{2} = \tfrac{1}{N}\sum_{i,j}(x_{ij} - \bar{x})^{2}$.
LISA classification yields four spatial cluster types:

\begin{table}[H]
\centering
\caption{LISA quadrant classification and fire-risk interpretation}
\label{tab:lisa}
\begin{tabular}{@{}lll@{}}
\toprule
\textbf{LISA Class} & \textbf{Description} & \textbf{Fire Risk Implication} \\
\midrule
High--High (HH) & High NDVI, high-NDVI neighbourhood &
    Dense contiguous forest; interior fire risk \\
Low--Low (LL)   & Low NDVI, low-NDVI neighbourhood &
    Degraded/sparse forest; low fuel load \\
High--Low (HL)  & High NDVI, low-NDVI neighbourhood &
    Forest fragment at WUI; elevated edge fire risk \\
Low--High (LH)  & Low NDVI, high-NDVI neighbourhood &
    Recent burn scar / clearcut within forest \\
\bottomrule
\end{tabular}
\end{table}

The LISA cluster class is incorporated as a \textbf{categorical feature}
in the ML pipeline, encoding spatial neighbourhood context that
pixel-level NDVI values alone cannot express.

% =============================================================
\subsection{NDVI--Fire Threshold Detection via Breakpoint Analysis}
\label{sec:threshold}
% =============================================================

\begin{novelbox}{Novel Contribution~5: Non-linear NDVI--Fire Relationship}
\citet{biswas2025} implicitly assumed a monotonic linear relationship
between NDVI and fire probability within the MaxEnt framework. This
study explicitly models the \textbf{non-linearity} of the NDVI--fire
relationship by detecting the critical threshold $\theta^{*}$ below
which fire occurrence probability increases sharply.
\end{novelbox}

\subsubsection{Piecewise Logistic Fire--NDVI Model}
\label{sec:piecewise}

Let $P(F = 1 \mid \text{NDVI} = x)$ denote the probability of fire
occurrence given pixel NDVI value $x$. A piecewise logistic model with
a single structural breakpoint $\theta$ is:

\begin{equation}
    P(F = 1 \mid x) =
    \begin{cases}
        \sigma\!\left(\alpha_{1} + \beta_{1}\,x\right) &
            \text{if } x \leq \theta \\[6pt]
        \sigma\!\left(\alpha_{2} + \beta_{2}\,x\right) &
            \text{if } x > \theta
    \end{cases}
    \label{eq:piecewise_logistic}
\end{equation}

where $\sigma(z) = (1 + e^{-z})^{-1}$ is the logistic sigmoid.
The parameter vector $\bm{\psi} = \{\alpha_1, \beta_1, \alpha_2,
\beta_2, \theta\}$ is estimated by minimising the binary
cross-entropy (negative log-likelihood):

\begin{equation}
    \hat{\bm{\psi}} =
    \underset{\bm{\psi}}{\arg\min}\;
    \mathcal{L}(\bm{\psi}) =
    -\sum_{n=1}^{N}
    \Bigl[
        F_n \log P_n + (1 - F_n)\log(1 - P_n)
    \Bigr]
    \label{eq:nll}
\end{equation}

The breakpoint $\theta^{*}$ is found via grid search over the
10th--90th percentile of observed NDVI followed by gradient-based
fine-tuning (Nelder--Mead optimisation). The estimated $\theta^{*}$
represents the \textbf{critical vegetation density threshold for fire
ignition} in each biogeographic zone.

\subsubsection{Biogeographic Zone-Stratified Thresholds}
\label{sec:zone_thresh}

India's forest types span five major biogeographic zones
\citep{rodgers1988}: (1)~Northeastern tropical, (2)~Western Ghats
tropical wet, (3)~Deccan dry deciduous, (4)~Central Indian mixed, and
(5)~Himalayan sub-alpine. A single national $\theta^{*}$ is
ecologically inadequate. Zone-stratified thresholds $\{\theta^{*}_{z}\}_{z=1}^{5}$
are estimated independently and incorporated as zone-specific NDVI
risk indicators in the final susceptibility mapping.

% =============================================================
\subsection{NDVI Feature Engineering Summary}
\label{sec:ndvi_features}
% =============================================================

Table~\ref{tab:ndvi_features} summarises all NDVI-derived features
entering the ML pipeline, expanding from the single raw NDVI layer
of \citet{biswas2025} to a nine-feature NDVI representation.

\begin{table}[H]
\centering
\caption{Complete NDVI-derived feature set for forest fire
         susceptibility mapping}
\label{tab:ndvi_features}
\begin{tabular}{@{}clllc@{}}
\toprule
\textbf{ID} & \textbf{Feature} & \textbf{Equation} &
\textbf{Physical Meaning} & \textbf{Novel?} \\
\midrule
F1 & QA-filtered NDVI $\overline{\text{NDVI}}_{\text{QA}}$
   & Eq.~\eqref{eq:qa_filter}
   & Mean vegetation density (cloud-screened)
   & Partial \\
F2 & Climatological mean $\bar{\mu}^{(m)}$
   & Eq.~\eqref{eq:clim_mean}
   & Long-term seasonal baseline
   & No \\
F3 & Monthly anomaly $\delta$
   & Eq.~\eqref{eq:ndvi_anomaly}
   & Inter-annual vegetation stress
   & \textbf{Yes} \\
F4 & STL trend $T$
   & Eq.~\eqref{eq:stl}
   & 22-year vegetation change trajectory
   & \textbf{Yes} \\
F5 & STL residual $R$
   & Eq.~\eqref{eq:stl}
   & Anomalous events (drought, fire scar)
   & \textbf{Yes} \\
F6 & Mann--Kendall $\tau$
   & Eq.~\eqref{eq:mk_tau}
   & Browning/greening trend significance
   & \textbf{Yes} \\
F7 & CVSI$(k^{*})$
   & Eq.~\eqref{eq:cvsi}
   & Pre-fire fuel moisture deficit
   & \textbf{Yes} \\
F8 & Local Moran's $I_{ij}$ / LISA
   & Eq.~\eqref{eq:local_morans_i}
   & Spatial neighbourhood vegetation cluster
   & \textbf{Yes} \\
F9 & $\mathbb{1}[\text{NDVI} < \theta^{*}_{z}]$
   & Eq.~\eqref{eq:piecewise_logistic}
   & Critical ignition threshold indicator
   & \textbf{Yes} \\
\midrule
\multicolumn{4}{l}{\textit{Biswas et al.\ (2025): 1 NDVI feature}} &
    --- \\
\multicolumn{4}{l}{\textit{This work: 9 NDVI-derived features}} &
    7 novel \\
\bottomrule
\end{tabular}
\end{table}

% =============================================================
\subsection{Integration with MaxEnt and ML Models}
\label{sec:integration}
% =============================================================

The nine NDVI-derived features (Table~\ref{tab:ndvi_features}) are
integrated into the prediction pipeline alongside the 14 other
variables of \citet{biswas2025} (Table~2 therein), yielding a
\textbf{23-dimensional feature space} compared to the 15-dimensional
space of the baseline study. In the MaxEnt framework, each NDVI
feature enters as a separate continuous environmental layer. In the
supervised ML pipeline (Random Forest, XGBoost, SVM), SHAP
(SHapley Additive exPlanations) analysis \citep{lundberg2017} is
applied to decompose each feature's marginal contribution — providing
model-agnostic attribution that extends beyond the permutation
importance reported in \citet{biswas2025}.

% =============================================================
% BIBLIOGRAPHY
% =============================================================
\newpage
\bibliographystyle{abbrvnat}
\bibliography{ndvi_references}

% ── Inline bibliography (fallback if .bib file not available) ──
% Remove the \bibliography{} line above and uncomment below:
%
% \begin{thebibliography}{99}
%
% \bibitem[Biswas et al.(2025)]{biswas2025}
% Biswas, U., Mahato, S., and Joshi, P.K. (2025).
% \newblock Forest fire occurrence probability modelling using MaxEnt in India.
% \newblock {\em Environmental Science and Pollution Research}, 32:4856--4878.
%
% \bibitem[Didan(2021)]{didan2021}
% Didan, K. (2021).
% \newblock {MODIS/Terra Vegetation Indices Monthly L3 Global 0.05Deg CMG V061}.
% \newblock NASA EOSDIS Land Processes DAAC.
% \newblock \url{https://doi.org/10.5067/MODIS/MOD13C2.061}
%
% \bibitem[Cleveland et al.(1990)]{cleveland1990stl}
% Cleveland, R.B., Cleveland, W.S., McRae, J.E., and Terpenning, I. (1990).
% \newblock {STL}: A seasonal-trend decomposition procedure based on loess.
% \newblock {\em Journal of Official Statistics}, 6(1):3--73.
%
% \bibitem[Mann(1945)]{mann1945}
% Mann, H.B. (1945).
% \newblock Nonparametric tests against trend.
% \newblock {\em Econometrica}, 13(3):245--259.
%
% \bibitem[Moran(1950)]{moran1950}
% Moran, P.A.P. (1950).
% \newblock Notes on continuous stochastic phenomena.
% \newblock {\em Biometrika}, 37:17--23.
%
% \bibitem[Chuvieco et al.(2004)]{chuvieco2004}
% Chuvieco, E., Cocero, D., Riano, D., Martin, P., Aguado, I.,
% de la Riva, J., and Burgan, R. (2004).
% \newblock Combining NDVI and surface temperature for the estimation of
% live fuel moisture content in forest fire danger rating.
% \newblock {\em Remote Sensing of Environment}, 92:322--331.
%
% \bibitem[Flannigan et al.(2000)]{flannigan2000}
% Flannigan, M.D., Stocks, B.J., and Wotton, B.M. (2000).
% \newblock Climate change and forest fires.
% \newblock {\em Science of the Total Environment}, 262:221--229.
%
% \bibitem[Lundberg and Lee(2017)]{lundberg2017}
% Lundberg, S.M. and Lee, S.I. (2017).
% \newblock A unified approach to interpreting model predictions.
% \newblock {\em Advances in Neural Information Processing Systems}, 30.
%
% \bibitem[Rodgers and Panwar(1988)]{rodgers1988}
% Rodgers, W.A. and Panwar, H.S. (1988).
% \newblock {\em Planning a Protected Area Network in India}.
% \newblock Wildlife Institute of India, Dehradun.
%
% \end{thebibliography}

\end{document}
